\section{Algoritma dan Simulasi}
Langkah perhitungan dapat dituliskan sebagai algoritma, seperti pada \autoref{alg:contoh}.

\begin{algorithm}
  \caption{Contoh algoritma}\label{alg:contoh}
  \begin{algorithmic}[1]
    \Require $x_0\in\R$, toleransi $\varepsilon>0$
    \Ensure Hampiran $x^*$
    \State $k\gets 0$
    \Repeat
      \State $x_{k+1}\gets T(x_k)$
      \State $k\gets k+1$
    \Until{$|x_{k}-x_{k-1}|<\varepsilon$}
    \State \Return $x_k$
  \end{algorithmic}
\end{algorithm}

Hasil simulasi disajikan pada \ref{tab:hasil}. Judul tabel diletakkan di atas tabel, sedangkan judul gambar di bawah gambar.

\begin{table}[H]
  \centering
  \caption{Contoh tabel hasil simulasi}
  \label{tab:hasil}
  \renewcommand{\arraystretch}{1.2}
  \begin{tabular}{|c|c|c|}
    \hline
    \textbf{Iterasi} & \textbf{$x_k$} & \textbf{Galat} \\ \hline
    1 & 0,500 & 0,500 \\ \hline
    2 & 0,750 & 0,250 \\ \hline
    3 & 0,875 & 0,125 \\ \hline
  \end{tabular}
\end{table}

\begin{figure}[H]
  \centering
  \includegraphics[width=0.3\textwidth]{lambang-its}
  \caption{Contoh gambar (sumber: dokumen pribadi)}
  \label{fig:contoh}
\end{figure}
